\section*{第 \ref{ch:b1:stratigraphic-scale} 章——地层年代表}

\begin{solution}{exo:b1:stratigraphic-scale:1}
(a)~\omterm{def:b1:stratigraphic-scale:chronostratigraphy}{年代地层单位}，系。(b)~\omterm{def:b1:stratigraphic-scale:lithostratigraphy}{岩石地层单位}，组。
(c)~\omterm{def:b1:stratigraphic-scale:chronostratigraphy}{地质年代单位}，世（其岩石为下白垩统）。
(d)~生物地层单位，\omterm{def:b1:stratigraphic-scale:biozone}{生物带}。(e)~\omterm{def:b1:stratigraphic-scale:chronostratigraphy}{年代地层单位}，阶。
(f)~\omterm{def:b1:stratigraphic-scale:chronostratigraphy}{地质年代单位}，世。(g)~\omterm{def:b1:stratigraphic-scale:lithostratigraphy}{岩石地层单位}，段。
\end{solution}

\begin{solution}{exo:b1:stratigraphic-scale:2}
古生代 $538.8 - 251.902 = \qty{286.9}{\Ma}$，占显生宙的 \qty{53.2}{\%}；
中生代 $251.902 - 66.00 = \qty{185.9}{\Ma}$，\qty{34.5}{\%}；
新生代 \qty{66.0}{\Ma}，\qty{12.2}{\%}。显生宙占
自冥古宙底界以来时间的 $538.8/4567 = \qty{11.8}{\%}$。
\end{solution}

\begin{solution}{exo:b1:stratigraphic-scale:3}
A、B、C、D 依此顺序沉积（\omterm{def:b1:stratigraphic-scale:principles}{叠加原理}），且为水平沉积（\omterm{def:b1:stratigraphic-scale:principles}{原始水平律}），
故其倾斜为后来的事件。岩墙 E 切穿 A 至 D，
故比 D 年轻（\omterm{def:b1:stratigraphic-scale:principles}{切割律}）；E 侵入与倾斜的先后，
由所给描述无法判定。侵蚀面切穿 D 与 E，
故比二者都年轻（\omterm{def:b1:stratigraphic-scale:principles}{切割律}）。F 与 G 随后沉积其上，
（\omterm{def:b1:stratigraphic-scale:principles}{叠加原理}）。断层错断 A 至 F 而不错断 G，
故其活动在 F 之后、G 之前（\omterm{def:b1:stratigraphic-scale:principles}{切割律}）。顺序为：A、B、C、D；倾斜与 E 的侵入（先后未定）；
侵蚀；F；断层作用；G。
\end{solution}

\begin{solution}{exo:b1:stratigraphic-scale:4}
$r = 30/(85.40 - 84.60) = \qty{37.5}{m}$ 每百万年。首现
位于下部火山灰之后 $12/37.5 = \qty{0.32}{\Ma}$ 处，即 $85.40 - 0.32 =
\qty{85.08}{\Ma}$。
\end{solution}

\begin{solution}{exo:b1:stratigraphic-scale:5}
(a)~以高度为纵轴画出四根竖条。(b)~B 与 C 同时出现于
\qtyrange{26}{50}{m} 之间，为共存\omterm{def:b1:stratigraphic-scale:biozone}{延限带}，厚 \qty{24}{m}。\omterm{def:b1:stratigraphic-scale:biozone}{间隔
带}自 B 的首现至 D 的首现，跨越 \qtyrange{12}{41}{m}，
\qty{29}{m}。(c)~在第二剖面中可确定 B 的首现（\qty{3}{m}）
与 C 的首现（\qty{9}{m}），因此可识别 (b) 中\omterm{def:b1:stratigraphic-scale:biozone}{间隔带}的底界与
共存\omterm{def:b1:stratigraphic-scale:biozone}{延限带}的底界均可识别；其顶界（D 的首现）及
在剖面~1 中所见的重叠带顶界均无法核对，因 D 缺失。
B 与 C 的首现（$\qty{12}{m} \leftrightarrow
\qty{3}{m}$，$\qty{26}{m} \leftrightarrow \qty{9}{m}$）最为可靠；
末现（B 位于 \qty{50}{m} 与 \qty{15}{m}）更多取决于保存状况。
\end{solution}

\begin{solution}{exo:b1:stratigraphic-scale:6}
(a)~$60/20 = \qty{3}{\Ma}$：陆向井中砂岩的顶部
年轻 300 万年。(b)~100 万年的泥质沉积，速率为 \qty{30}{m}
每百万年：火山灰位于砂岩顶部之上 \qty{30}{m} 处。
(c)~火山灰降落时，海岸线已较靠海井向陆移动 \qty{20}{km}，
故仍在陆向井以海侧 \qty{40}{km} 处，而陆向井当时位于陆地上。
火山灰落于该处陆面上，其上为 200 万年后海水沉积的砂岩，
除非其先被冲走——陆上松散层很可能如此。
同一火山灰层在一口井中位于组界之上，
在另一口井中则位于组界之下：该界线穿越了时间。
\end{solution}

\begin{solution}{exo:b1:stratigraphic-scale:7}
(a)~洛赫科夫阶 $1.36/419.62 = \qty{0.32}{\%}$；印度阶 $0.024/251.902 =
\qty{0.0095}{\%}$，约小 34 倍。印度阶底界位于
经精确测年的 GSSP 剖面本身的火山灰层之间；较大的不确定度表明，
该界线年龄须由较远处的定年层位推算。
(b)~$251.941 - 251.880 = \qty{0.061}{\Ma} = 61$ 千年，其
不确定度为 $\sqrt{37^2 + 31^2} = 48$ 千年，故已发表的数值为
\num[uncertainty-mode=separate]{60 \pm 48} 千年。
\end{solution}

\begin{solution}{exo:b1:stratigraphic-scale:8}
(a)~$m = 3000 \times \frac43\pi\,(5000)^3 = \qty{1.57e15}{kg}$；尘埃
质量 $60m = \qty{9.42e16}{kg}$。(b)~地球表面积为 $4\pi(6.378\times10^6)^2
= \qty{5.11e14}{m^2}$：每平方米 \qty{184}{kg}。(c)~$184/2500 =
\qty{0.074}{m}$，约 \qty{7}{cm}。(d)~厚度与直径的
立方成正比：直径 \qty{6}{km} 时为 \qty{1.6}{cm}，\qty{14}{km} 时为 \qty{20}{cm}。数量级
即厘米至分米，与边界黏土的相当，
在 El Kef 厚 \qty{50}{cm}；比较只能粗略，因为尘埃
不可能均匀降落，且黏土层还汇入了同期
的普通沉积物。
\end{solution}

\begin{solution}{exo:b1:stratigraphic-scale:9}
(a)~$\num{60e3} \times \num{2e-6} = \qty{0.12}{m}$，\qty{12}{cm}；取 12 千年与
108 千年，分别为 \qty{2.4}{cm} 与 \qty{21.6}{cm}。(b)~$\num{32e3} \times
\num{10e-6} = \qty{0.32}{m}$。(c)~灭绝事件占据厘米至分米量级的岩层，
故样品须逐厘米采集；压缩段中单个样品
厚 \qty{10}{cm}，平均跨越 5 万年——
占整个事件的大部分。
\end{solution}

\begin{solution}{exo:b1:stratigraphic-scale:10}
(a)~$u = R/s = 10 \times 2000 = \qty{20000}{m}$，即每百万年
\qty{20}{km}；在 \qty{30}{km} 范围内，$\Delta t = \qty{1.5}{\Ma}$。
(b)~平原如今的地表为 $z = s x + a t$。海岸线满足
$s x + a t = z_0 + R t$，故 $x(t) = (z_0 + (R - a)t)/s$，且
$\mathrm dx/\mathrm dt = (R - a)/s$，当 $a < R$ 时为正（向陆）。当
$a > R$ 时速度为负：海岸线向海移动，即海退，且
海岸线砂岩此时向海一侧更年轻。(c)~一段长 \qty{200}{m} 的悬崖，
以 \qty{20}{km} 每百万年的速度计，$0.2/20 = \qty{0.01}{\Ma}$——即 1 万年
的穿时，远低于化石或多数年代测定所能分辨的程度；接触面两侧的
岩石沿整个接触面处处相同。
\end{solution}

\begin{solution}{exo:b1:stratigraphic-scale:11}
(a)~$f = (25 - 10)/(50 - 10) = 0.375$，且 $t = 100.50 - 0.375 \times 0.80 =
\qty{100.20}{\Ma}$，其中
\[
  \sigma_t = \sqrt{(0.625 \times 0.05)^2 + (0.375 \times 0.03)^2} = \qty{0.033}{\Ma}.
\] (b)~$g(f) = (1 - f)^2\sigma_1^2 + f^2\sigma_2^2$
有 $g'(f) = -2(1 - f)\sigma_1^2 + 2f\sigma_2^2$，其零点为 $f^* =
\sigma_1^2/(\sigma_1^2 + \sigma_2^2) = 0.0025/0.0034 = 0.735$，更靠近定年较精确的
层位；此时 $\sigma_t = \sigma_1\sigma_2/\sqrt{\sigma_1^2 + \sigma_2^2} =
\qty{0.026}{\Ma}$，小于两者中的任一个。(c)~$g$ 为开口向上的抛物线，故
在 $[0, 1]$ 上，其最大值在端点取得：$g(0) = \sigma_1^2$，$g(1) = \sigma_2^2$。
因此两层位之间，$\sigma_t \le \max(\sigma_1, \sigma_2)$。
\end{solution}

\begin{solution}{exo:b1:stratigraphic-scale:12}
(a)~边界位于深度的一半处：$f = 0.5$，\qty{29.50}{\Ma}。(b)~
取 $\varphi_0\lambda = 0.63 \times 1960 = \qty{1234.8}{m}$，则
\qtyrange{600}{1000}{m} 段的固体厚度为 $400 - 1234.8\,(e^{-600/1960} - e^{-1000/1960}) =
\qty{232.2}{m}$，\qtyrange{200}{600}{m} 段为 \qty{194.2}{m}；合计
\qty{426.3}{m}。下部火山灰与边界之间沉积的颗粒所占比例为
$232.2/426.3 = 0.545$，年龄则为 $30.0 - 0.545 =
\qty{29.46}{\Ma}$。(c)~第二个估计值年轻 \num{45} 千年。
它更为可靠：较深的一半失去的水更多，因此
每米所含颗粒和所代表的时间都比浅部一半更多（较深一半的平均孔隙度
为 \num{0.42}，浅部一半为 \num{0.51}）。
\end{solution}

\subsection*{问题：在盆地中对比三个剖面}

\begin{solution}{pb:b1:stratigraphic-scale:1}
\begin{enumerate}
  \item 三根柱状图可取 \qty{1}{cm} 代表 \qty{10}{m} 绘制，并由火山灰层
        与带界线相连；凯尔砂岩楔状体由 A 处的
        \qty{45}{m} 变薄至 C 处的 \qty{6}{m}。
  \item 带 2：A $65.8 - 36.6 = \qty{29.2}{m}$，B \qty{20.8}{m}，C
        \qty{10.0}{m}。带 3 的记录厚度：A \qty{14.2}{m}，B \qty{14.0}{m}，C
        \qty{7.2}{m}。
  \item 带 1 始于 \emph{V.~primum} 的首现，在每个剖面的底部都可找到它：
        其首现位于各剖面之下。
  \item 它只见于砂岩中：它随岩相变化，而
        岩相在各地的时代并不相同。在 A 处，它出现于带 1
        和带 2 中，在 B 与 C 处仅出现于带 1 中。
  \item 三处均出现在带 2 中（A \qty{49.0}{m}，B \qty{40.0}{m}，C
        \qty{19.0}{m}）。它在此三处彼此一致，因此
        可用来细分带 2，但末现取决于保存情况；
        应先在更多剖面中加以检验。
  \item 各连线分别经过 \qty{36.6}{m}、\qty{31.2}{m} 和 \qty{14.8}{m}，再经过
        \qty{65.8}{m}、\qty{52.0}{m} 和 \qty{24.8}{m}。带 2 的厚度之比为
        $29.2 : 20.8 : 10.0$，约为 $2.9 : 2.1 : 1$：若该带在各处历时相同，
        则沉积物在 A 处的堆积速度约为
        C 处的三倍。
  \item A：带 2（$45.0 > 36.6$）；B：带 1（$20.0 < 31.2$）；C：带 1
        （$6.0 < 14.8$）。
  \item 否：它在 A 处位于带 2，在 B 与 C 处位于带 1，因此越往西越年轻。
        海相页岩向西扩展，覆盖在滨岸砂岩之上：
        海岸线向西、向内陆推移，形成海侵。
  \item 火山灰在极短的地质瞬间落满整个盆地，并沉落于
        任何沉积物之上；砂岩顶面随移动的海岸线变化，
        因而是\omterm{def:b1:stratigraphic-scale:diachronous}{穿时的}。
  \item \omterm{def:b1:stratigraphic-scale:principles}{横向连续律}（该层跨越整个盆地延伸），以及
        \omterm{def:b1:stratigraphic-scale:principles}{叠加原理}（每个剖面中火山灰~1 都位于火山灰~2 之下，这可检验
        对比结果）。
  \item A 处 $(71.0 - 5.0)/1.20 = \qty{55.0}{m}$ 每百万年；B 处
        $48.0/1.20 = \qty{40.0}{}$；C 处 $24.0/1.20 = \qty{20.0}{}$。
  \item A 处最靠近陆地，那里是沉积物的供给地；
        进入东侧较深水域的沉积物更少。
  \item $0.1/55 = \num{1.8}$ 千年（A 处），$\num{2.5}$ 千年（B 处），$\num{5.0}$
        千年（C 处）：A 处对时间的分辨最为精细。
  \item A 处 $152.40 - 40.0/55.0 = \qty{151.67}{\Ma}$；B 处 $152.40 - 12.0/40.0 =
        \qty{152.10}{\Ma}$；C 处 $152.40 - 3.0/20.0 = \qty{152.25}{\Ma}$。
  \item 由 C 至 B：$20/(152.25 - 152.10) = \qty{133}{km}$ 每百万年；由 B 至 A：
        $25/(152.10 - 151.67) = \qty{58.5}{km}$ 每百万年。速度
        减慢了。
  \item A 处 $152.40 - 60.8/55.0 = \qty{151.29}{\Ma}$；B 处 $152.40 - 44.0/40.0 =
        \qty{151.30}{\Ma}$；C 处 $152.40 - 21.8/20.0 = \qty{151.31}{\Ma}$。
  \item 它们跨越 \num{16} 千年，小于每个年龄的不确定度
        （约 \num{28} 千年，见问题~19）：首现在
        测年分辨率范围内于整个盆地同步出现，
        这正是它成为良好\omterm{def:b1:stratigraphic-scale:biozone}{标准化石}的原因。
  \item A 处 $152.40 - 31.6/55.0 = \qty{151.826}{\Ma}$；B 处 $152.40 - 23.2/40.0 =
        \qty{151.820}{\Ma}$；C 处 $152.40 - 11.8/20.0 = \qty{151.810}{\Ma}$。
  \item $f = 0.479$、$0.483$、$0.492$；$\sigma_t = \sqrt{(1 - f)^2 + f^2} \times
        0.04 = \qty{0.028}{\Ma}$，三处均为 \num{28} 千年。
  \item 平均 \qty{151.819}{\Ma}；离散范围 $151.826 - 151.810 = \num{16}$ 千年，
        小于每个年龄的不确定度：三者一致。
  \item 插值假定从火山灰~1 到火山灰~2 之间速率一致。在 A 处，该间隔
        内含至 \qty{45}{m} 的滨岸砂岩，其上为页岩；
        砂和泥未必以相同速率沉积，且二者的
        \omterm{def:b1:diagenesis-basins:compaction}{压实}程度不同（见\cref{ch:b1:diagenesis-basins}）。若砂的沉积
        较快，则砂岩段所含的时间少于其厚度所暗示，
        首现年龄便比计算值更老，
——A 处的年龄确实是三者中最老的。
  \item B 处：首现位于海相页岩中，比砂岩顶面高 \qty{11}{m}，
        该层位附近岩相无变化，按该速率
        得每 \qty{10}{cm} \num{2.5} 千年。A 处对时间的分辨更好，但该层位
        落在滨岸砂岩中，低于一处岩相变化 \qty{8.4}{m}；
        C 处为压缩层段，每
        \qty{10}{cm} \num{5} 千年。就分辨率而言 C 处最差，就岩相而言 A 处最差。
  \item 由两个首现的平均年龄得 $151.819 - 151.302 =
        \qty{0.52}{\Ma}$：约五十万年。
  \item 该边界为 \qty[uncertainty-mode=separate]{151.82(3)}{\Ma}。
\end{enumerate}
\end{solution}
